\section{Branch load matters, but branch-local shortcut is unmeasured (RQ4)}
\label{sec:rq4}

\begin{figure}[t]
\centering
\includegraphics[width=\linewidth]{fig4_branch_blocked.pdf}
\caption{Branch-locality in the FBCSP-LGG fusion backbone. The spatial branch is load-bearing (largest
ablation drop and highest fusion-gate weight on both datasets); graph/temporal branches are
neutral-to-starved. Per-branch \emph{leakage} (L1) and \emph{reliance} (L5) are not measured, so RQ4 is
blocked (no frozen checkpoint exists to measure them without retraining).}
\label{fig:branch}
\end{figure}

A multi-branch backbone lets us ask whether leakage \emph{means} something different in a load-bearing
branch than in a non-load-bearing one. In FBCSP-LGG the spatial branch is clearly load-bearing:
zeroing it drops task-bAcc by $7.4$ points on 2a and $8.8$ on 2015, and it receives the largest
fusion-gate weight ($0.49$/$0.57$); the graph and temporal branches are neutral-to-slightly-harmful and
starved after fusion (Figure~\ref{fig:branch}).

\paragraph{Blocked, not failed.} Answering the branch-locality question quantitatively requires
per-branch \emph{leakage} (L1) and per-branch \emph{reliance} (L5). Neither exists: there are no frozen
per-branch latent dumps and no per-branch probe on disk, and --- decisively --- \emph{no trained
checkpoint exists at all}. The training pipeline persisted only scalar summaries; the best-epoch state
was held in memory and never serialized, and a read-only search of the compute store returned zero
matching checkpoints. Producing the metric would require retraining, which we do not do. RQ4 is
therefore reported as blocked: we can state that the spatial branch is load-bearing, but we
\emph{cannot} say ``spatial leakage is harmful,'' ``graph leakage is benign,'' or ``per-branch
conditional MI predicts reliance.'' The blocked status is itself a finding: a branch-local shortcut
claim requires an instrument that does not currently exist.
